% =====================================================================
% FaaS-MAGCAA: greedy coalition-auction allocation for FRALB.
% Meant to be \input{} (or pasted) into the paper. The notation recap
% below is removable when the host paper has already defined FRALB.
% =====================================================================

\section{FaaS-MAGCAA: greedy coalition auction allocation}
\label{sec:magcaa}

\textsc{FaaS-MAGCAA} adapts the task-selection and winner-consensus
mechanism of the Greedy Coalition Auction Algorithm (GCAA) of
\citet{braquetbakolas2021} to Function Replica Allocation and Load Balancing
(FRALB). It is included as an experimental decentralized baseline rather
than as a new auction-theoretic mechanism. At every control period, local
problems first determine the load that each node would like to offload and
the replicas already available. Repeated auction rounds then match unit
loads to neighbouring execution capacity by greedy utility ranking.

This baseline deliberately omits two mechanisms used by other methods in
the paper. First, its price array is fixed at zero: there is no adaptive
price update as in \textsc{FaaS-MADeA}. Second, its replica-seller array is
fixed at zero: there is no memory or replica market, and unmatched load
receives the null allocation. Consequently, the method provides neither the
price or dual certificate associated with \textsc{FaaS-MALD} nor the
equilibrium certificate associated with \textsc{FaaS-MAPG}. Its role is to
isolate the empirical behaviour of a greedy GCAA-style selection and
consensus rule inside the common FRALB execution pipeline.
Classical price-based auction algorithms and distributed numerical methods
provide broader background for such coordination
\citep{bertsekas1988,bertsekas1989}; FaaS-MAGCAA's fixed zero price means that
it does not inherit their price-adjustment results or guarantees.

% =====================================================================
% BEGIN self-contained notation recap.
% This subsection restates only the FRALB notation needed by FaaS-MAGCAA.
% REMOVE this whole fenced block when inserting the section into a paper
% that has already introduced these symbols and the capacity model.
% =====================================================================
\subsection{Notation and capacity model}
\label{sec:magcaa-notation}

Let $\mathcal N$ be the set of compute nodes, $\mathcal F$ the function set,
and $N_i\subseteq\mathcal N$ the one-hop neighbourhood of node $i$. At a
control period, node $i$ receives workload $\lambda_i^f$ for function $f$.
Its local solve partitions this workload into locally executed load $x_i^f$,
load $\omega_i^f$ offered to the coordination mechanism, and rejected or
Cloud-forwarded load $z_i^f$. The coordination variable $y_{ij}^f$ is the
amount originating at $i$ and executed by neighbour $j$. Node $j$ runs
$r_j^f$ replicas, each governed by maximum utilization $U_{\max}^f$ and by
the per-request service demand $D_j^f$ used by the implementation.

Before auction round $h$, the implementation computes execution capacity
and residual capacity as
\begin{equation}
  \operatorname{Cap}_j^f(h)=
    \frac{r_j^f(h)U_{\max}^f}{D_j^f},
  \qquad
  C_j^f(h)=\max\!\left\{0,
    \frac{r_j^f(h)U_{\max}^f}{D_j^f}
    -x_j^f(h)-\sum_{i\in\mathcal N}y_{ij}^f(h)
  \right\}.
  \label{eq:magcaa-residual-capacity}
\end{equation}
Thus $C_j^f$ is measured in load units, not replicas. The bid generator is
given the looser blackboard quantity
$B_j^f=\max\{0,\operatorname{Cap}_j^f-x_j^f\}$, whereas final acceptance
uses $C_j^f$, which also subtracts inbound assignments accumulated in $y$.

For a buyer $i$, seller $j$, and function $f$, the inherited helper computes
$\beta_{ij}^f-p_j^f-w_L\ell_{ij}-w_Fq_i^f(h)$. FaaS-MAGCAA fixes
$p_j^f=0$, so its specialized baseline utility is
\begin{equation}
  u_{ij}^f(h)=\beta_{ij}^f-w_L\ell_{ij}-w_Fq_i^f(h),
  \label{eq:magcaa-utility}
\end{equation}
where $\beta_{ij}^f$ is the FRALB offloading benefit, $\ell_{ij}$ is network
latency, and $q_i^f$ is a within-period fairness counter. The configurations
used for this baseline set $w_L=w_F=0$, so
$u_{ij}^f=\beta_{ij}^f$. The symbol $\ell_{ij}$ in
Eq.~\eqref{eq:magcaa-utility} denotes latency and is distinct from the
served-load quantity $x_j^f+\sum_i y_{ij}^f$ computed internally by the
capacity helper.

The no-ping-pong state for function $f$ is represented by
\begin{equation}
  S_i^f(y)=\mathbf 1\!\left\{\sum_j y_{ij}^f>0\right\},
  \qquad
  R_i^f(y)=\mathbf 1\!\left\{\sum_k y_{ki}^f>0\right\},
  \label{eq:magcaa-send-receive}
\end{equation}
where $S_i^f$ and $R_i^f$ mean, respectively, that $i$ is sending or
receiving function-$f$ load. Feasibility requires that they are not both
true for the same node and function.

\begin{table}[t]
\centering
\caption{Notation used in the FaaS-MAGCAA section.}
\label{tab:magcaa-notation}
\small
\begin{tabular}{@{}l l p{0.59\linewidth}@{}}
\toprule
Nodes, functions & $\mathcal N,\mathcal F$ & compute nodes and deployable functions.\\
Neighbourhood & $N_i$ & sellers visible to buyer node $i$.\\
Input workload & $\lambda_i^f$ & function-$f$ load arriving at node $i$.\\
Local load & $x_i^f$ & load executed at its origin node $i$.\\
Auction load & $\omega_i^f$ & residual load of active buyer--function pair $(i,f)$.\\
Routed load & $y_{ij}^f$ & load sent from buyer $i$ to seller $j$.\\
Cloud/rejected load & $z_i^f$ & load not executed locally or by a neighbour.\\
Replicas & $r_i^f$ & number of function-$f$ replicas at node $i$.\\
Residual capacity & $C_j^f$ & seller capacity after local $x_j^f$ and inbound $\sum_i y_{ij}^f$.\\
Benefit, latency & $\beta_{ij}^f,\ell_{ij}$ & offloading benefit and network-latency input.\\
Fairness state & $q_i^f$ & number of prior inner rounds in the period in which $(i,f)$ has an allocation.\\
Utility & $u_{ij}^f$ & consensus ranking value in Eq.~\eqref{eq:magcaa-utility}.\\
Sending, receiving & $S_i^f,R_i^f$ & predicates in Eq.~\eqref{eq:magcaa-send-receive}.\\
Iteration limits & $H_{\max},T_{\max}$ & maximum auction rounds and accumulated solver-time limit.\\
\bottomrule
\end{tabular}
\end{table}
% ===================== END self-contained notation recap =====================

\subsection{From GCAA tasks to FRALB unit allocations}
\label{sec:magcaa-mapping}

In the original GCAA, an agent selects at most one task during an auction
iteration, agents proposing the same task are compared by utility, and the
winner is finalized. Several agents may ultimately form a coalition on the
same task over successive iterations; the null task remains available when
no task yields an acceptable assignment \citep{braquetbakolas2021}. In the
FRALB adaptation, every active buyer--function pair $(i,f)$ with
$\omega_i^f>0$ plays the role of one GCAA agent. A neighbouring
seller--function pair $(j,f)$ with advertised capacity plays the role of a
task. The translation is deliberately unit based: the runner rejects any
configuration in which \texttt{unit\_bids} is not true, and every bid row
therefore has demand $d=1$.

This mapping changes the meaning of coalition formation. A seller--function
task accepts at most one buyer in one FaaS-MAGCAA round, but its capacity may
be filled by buyers accepted in later rounds. The accumulated tensor $y$
therefore records the analogue of a coalition across rounds. Conversely, a
buyer--function pair is not permanently finalized after its first accepted
unit: if residual load remains, it can bid again in a later round. If no row
survives neighbourhood, convenience, capacity, and no-ping-pong checks, the
pair receives no allocation in that round. This is the FRALB counterpart of
GCAA's null assignment, without creating a replica to manufacture a task.

\paragraph{Utility and bid columns.}
The inherited bid generator produces both a utility column
$u_{ij}^f$ and a ranking-price column
\begin{equation}
  b_{ij}^f=p_j^f+\epsilon+\delta_{ij}^f,
  \label{eq:magcaa-b-column}
\end{equation}
where $\delta_{ij}^f$ is the gap to the next seller in the generator's
utility ordering. FaaS-MAGCAA fixes $p_j^f=0$ and does not call the adaptive
price-update routine. More importantly, its proposal and consensus code
ranks the \texttt{utility} column, not $b_{ij}^f$. Hence $\epsilon$ and
$\delta_{ij}^f$ may change the unused $b$ column but cannot change the
winner ranking; under the baseline configurations, that ranking is simply
descending $\beta_{ij}^f$.

\subsection{Proposal and consensus round}
\label{sec:magcaa-consensus}

Let $\mathcal B(h)$ denote the rows emitted by the bid generator in round
$h$. Because unit mode expands available quantities, $\mathcal B(h)$ can
contain repeated unit rows and rows for several sellers for the same
buyer--function pair. The proposal stage keeps exactly the row selected by
the data-frame operation
\begin{equation}
  \widehat b_i^f(h)\in
  \underset{b\in\mathcal B(h):\,(b.i,b.f)=(i,f)}{\operatorname{arg\,max}}
  b.\textit{utility}.
  \label{eq:magcaa-proposal}
\end{equation}
The implementation uses \texttt{groupby(...).idxmax()}, including its
data-frame semantics when equal maxima occur; no stronger universal
tie-breaking rule is part of the algorithmic contract.

The selected proposals are grouped by seller--function pair $(j,f)$. If
$S_j^f(y)=1$, the entire group is skipped because that seller already sends
the same function. Otherwise its contenders are visited in descending
utility. A contender $(i,f)$ is skipped if $R_i^f(y)=1$ or if
$C_j^f(h)<d$. The first contender passing both tests is accepted,
$y_{ij}^f$ is increased by $d=1$, and the round-local states are updated to
$S_i^f=1$ and $R_j^f=1$ before the next group is examined. The group then
terminates. A lower-ranked contender can therefore win when a higher-ranked
one is invalid under the receiving or capacity test; the code does not
blindly discard the whole task after rejecting its first row.

The state used to initialize these predicates is the cumulative
\texttt{current\_y} from all previous rounds. The immediate round-local
updates additionally prevent an allocation accepted earlier in the same
round from participating in a later ping-pong chain. Once all task groups
have been visited, the round increment is added to cumulative $y$.

\subsection{Implemented procedure}
\label{sec:magcaa-procedure}

Algorithm~\ref{alg:magcaa} states the complete per-period procedure. It
separates the greedy auction operation from the common FRALB machinery that
re-solves local models, validates feasibility, evaluates candidate welfare,
tracks two best candidates, checkpoints, and saves artifacts.

\begin{algorithm}[H]
\caption{\textsc{FaaS-MAGCAA} at one control period}
\label{alg:magcaa}
\begin{algorithmic}[1]
\Require period data, neighbourhoods, $H_{\max}$, $T_{\max}$, tolerance $\tau$
\State solve one local \textsc{LSP} for each node; obtain $x$, $r$, $\rho$, $U$, and initial components $\omega_i^{f,0}$
\State $y\gets0$; $\boldsymbol\omega\gets(\omega_i^{f,0})_{i,f}$; $q\gets0$; $p\gets0$; replica-seller vector $\widehat\rho\gets0$
\State require \texttt{unit\_bids}${}={}$true; set round $h\gets0$
\Repeat
  \State compute $\operatorname{Cap}$, $C$, and served load from $(x,y,r)$ by Eq.~\eqref{eq:magcaa-residual-capacity}
  \State $B\gets\max\{0,\operatorname{Cap}-x\}$
  \State $(\mathcal B,\mathcal B_{\rm mem},n_a)\gets\textsc{DefineBids}(\omega,B,p,\widehat\rho,q)$
  \Comment{$\mathcal B_{\rm mem}=\emptyset$ because $\widehat\rho=0$}
  \State $\widehat\omega\gets0$
  \If{$\mathcal B\ne\emptyset$}
    \State $\Delta y\gets0$; initialize $S,R$ from cumulative $y$
    \State keep one maximum-utility row per active $(i,f)$ as in Eq.~\eqref{eq:magcaa-proposal}
    \For{each resulting seller--function group $(j,f)$}
      \If{$S_j^f=0$}
        \For{each contender $(i,f,d)$ in descending utility}
          \If{$R_i^f=0$ and $C_j^f\ge d$}
            \State $\Delta y_{ij}^f\gets\Delta y_{ij}^f+d$; $S_i^f\gets1$; $R_j^f\gets1$
            \State \textbf{break}
          \EndIf
        \EndFor
      \EndIf
    \EndFor
    \State $y\gets y+\Delta y$; $\widehat\omega_i^f\gets\sum_j y_{ij}^f$
    \State increment $q_i^f$ for each $(i,f)$ with $\widehat\omega_i^f>0$
    \State run the pre-solve fixed-$y$ \textsc{LSPr} memory-feasibility check; abort with a runtime error if it fails
    \State solve fixed-$y$ \textsc{LSPr} models via \textsc{ComputeSocialWelfare}; update $x,r,\rho$
    \State $\omega_i^f\gets\omega_i^{f,0}-\widehat\omega_i^f$; set values within $\tau$ of zero to zero
  \EndIf
  \State combine $(x,y,r,\rho)$, tighten replicas to realized served load, and assert centralized feasibility
  \State update the best restricted-objective candidate and the best centralized-objective candidate
  \State $(\mathrm{stop},\mathrm{reason})\gets\textsc{CheckStoppingCriteria}(h,H_{\max},B,\omega,\widehat\omega,0,\mathcal B,\mathcal B_{\rm mem},\tau,\mathrm{runtime},T_{\max})$
  \If{not $\mathrm{stop}$} \State $h\gets h+1$ \EndIf
\Until{$\mathrm{stop}$}
\State decode the two tracked candidates and record the period objective and termination reason
\State checkpoint \textsc{LSP}/\textsc{LSPc} only if $t\bmod I_{\rm checkpoint}=0$ or $t=\texttt{max\_steps}-1$
\end{algorithmic}
\end{algorithm}

The zero in the stopping call is the default all-zero additional-replica
allocation created when the argument is passed as \texttt{None}; the
restricted allocation $\widehat\omega$ is cumulative, not merely the most
recent round increment. The variable $n_a$ is returned by the inherited bid
generator but is not used by this runner. Likewise, the local solution's
$\rho$ is retained for solution construction, while the separate all-zero
$\widehat\rho$ disables memory bids.

The runner always writes \texttt{config.json}. It writes \texttt{out.log}
only when \texttt{log\_on\_file=True}, and writes the sole plot artifact
\texttt{sp.png} only when plotting is enabled and both $N_f\le10$ and
$N_n\le10$. After all control periods, it joins the complete period solutions
and writes the final \textsc{LSP} and \textsc{LSPc} outputs. For each prefix
\texttt{LSP} and \texttt{LSPc}, these
are:
\begin{itemize}
  \item \texttt{<prefix>\_solution.csv};
  \item \texttt{<prefix>\_offloaded.csv};
  \item \texttt{<prefix>\_utilization.csv};
  \item \texttt{<prefix>\_replicas.csv};
  \item \texttt{<prefix>\_detailed\_fwd\_solution.csv}; and
  \item \texttt{<prefix>\_residual\_capacity.csv}.
\end{itemize}
The runner also writes
\texttt{obj.csv}, \texttt{termination\_condition.csv}, and
\texttt{runtime.csv}. When the checkpoint predicate in
Algorithm~\ref{alg:magcaa} holds, the component CSVs are written beneath
\texttt{LSP/<t>/} and \texttt{LSPc/<t>/}; those checkpoint directories are
not written for other periods.

\subsection{Round invariants}
\label{sec:magcaa-properties}

\begin{proposition}[At most one winner per task and round]
\label{prop:magcaa-one-winner}
For every round $h$ and seller--function pair $(j,f)$,
$\sum_i\Delta y_{ij}^f(h)$ contains at most one accepted unit.
\end{proposition}
\begin{proof}
All selected proposals for $(j,f)$ are processed in one data-frame group.
After the first valid contender is accepted, the inner contender loop
terminates. If no contender is valid, the group contributes zero. Because
the outer grouping visits $(j,f)$ only once, no second unit can be accepted
for that task during the same round.
\end{proof}

\begin{proposition}[Capacity at acceptance]
\label{prop:magcaa-capacity}
Every unit accepted for $(j,f)$ is no larger than the residual capacity
$C_j^f(h)$ computed at the start of that round.
\end{proposition}
\begin{proof}
The acceptance branch is entered only if $C_j^f(h)\ge d$, and mandatory
unit mode gives $d=1$. By Proposition~\ref{prop:magcaa-one-winner}, no other
unit for $(j,f)$ is accepted in that round, so the start-of-round residual
need not be decremented again within the group traversal. Before the next
round, $C_j^f$ is recomputed from the cumulative $y$ and the latest
restricted re-optimization. Thus each accepted unit respects the residual
capacity available at the time of its check, including load accumulated in
earlier rounds.
\end{proof}

\begin{proposition}[Inductive no-ping-pong preservation]
\label{prop:magcaa-no-ping-pong}
Starting from $y=0$, every cumulative allocation produced by
Algorithm~\ref{alg:magcaa} satisfies
$S_i^f(y)R_i^f(y)=0$ for every $(i,f)$.
\end{proposition}
\begin{proof}
The property is true at initialization. Assume it holds for cumulative $y$
at the beginning of a round. A proposed edge $i\to j$ is accepted only if
$R_i^f=0$ and $S_j^f=0$, so it cannot make an existing receiver send or an
existing sender receive. Immediately after acceptance, the algorithm sets
$S_i^f=1$ and $R_j^f=1$. Subsequent groups in the same round therefore
cannot use $i$ as a seller or $j$ as a buyer. Each accepted edge preserves
the invariant both relative to prior rounds and relative to earlier edges
in the current round. Induction over accepted edges and then over rounds
proves the claim.
\end{proof}

These propositions are local feasibility statements. They do not establish
global optimality, truthfulness, a market equilibrium, or a dual bound. In
particular, Proposition~\ref{prop:magcaa-capacity} concerns the residual at
each acceptance; overall FRALB feasibility is additionally checked by the
fixed-$y$ memory pre-check, the restricted local solves, and the centralized
feasibility assertion shown in Algorithm~\ref{alg:magcaa}.

\subsection{Stopping conditions and practical limits}
\label{sec:magcaa-stopping}

The runner reuses \texttt{check\_stopping\_criteria} from
\textsc{FaaS-MADeA} with its arguments specialized as shown in
Algorithm~\ref{alg:magcaa}. Its tests are ordered. It first stops when
$h\ge H_{\max}-1$. It next stops if all blackboard capacities $B$ are at
most $\tau$, or if all residual buyer loads $\omega$ are at most $\tau$.
The default all-zero additional-replica allocation then makes the inherited
``load cannot be assigned'' branch apply whenever cumulative
$\widehat\omega$ is also all zero. If that branch does not apply, the code
stops when both ordinary and memory bid tables are empty. Finally, it tests
whether accumulated solver runtime has reached $T_{\max}$. The later
``feasible solution found'' branch of the shared helper is inactive because
this call does not pass its optional \texttt{sp\_omega} and \texttt{sp\_y}
arguments.

The stopping helper is called only after bid resolution, any fixed-$y$
restricted re-solve, and candidate evaluation. Consequently, every entered
control period executes round $h=0$. For positive $H_{\max}$, the maximum
iteration guard bounds the number of completed rounds; it does not prevent
entry into the current round.

These guards cover exhaustion of offered load, exhaustion of advertised
capacity, absence of bidders, an initially unassignable allocation, the
iteration cap, and the time cap. They do not make every stalled state an
algorithmic fixed point: after some allocation has accumulated,
$\widehat\omega$ is nonzero, and bid rows rejected only by the stricter
residual-capacity or no-ping-pong checks may persist until another guard
fires. Reaching ``max iterations reached'' or the time limit therefore
returns a feasibility-checked candidate but is not a convergence
certificate.

The original GCAA proves completion within at most the number of agents
because each iteration finalizes an additional agent under that paper's
task-allocation protocol \citep{braquetbakolas2021}. That result does not
transfer automatically here. A FaaS-MAGCAA buyer--function pair may demand
several units, may bid again after an accepted unit, and competes under
capacity, re-optimization, eligibility, and no-ping-pong constraints.
Moreover, several distinct seller--function tasks can accept concurrently
in one round. Accordingly, this implementation relies on the practical
guards above and claims no inherited finite bound in terms of the number of
nodes or buyer--function pairs.

\subsection{Relation to GCAA and the other coordination baselines}
\label{sec:magcaa-positioning}

FaaS-MAGCAA retains the two recognizable operations of GCAA: an active agent
selects one currently best task, and contenders for a common task are
resolved greedily by utility. It does not reproduce the paper's distributed
bid-vector exchange or its permanent finalized-agent vector. Instead, the
shared runner materializes bid rows centrally for an experimental baseline,
resolves one winner per seller--function task and round, and carries the
accepted flow tensor across rounds. The original paper permits multiple
agents to form a task coalition through repeated finalizations; the FRALB
counterpart accumulates unit winners against execution capacity through
repeated fixed-$y$ re-optimization.

This construction also explains its narrower interpretation relative to
the paper's other baselines. Unlike \textsc{FaaS-MADeA}, it never updates
prices and never solicits replica memory. Unlike \textsc{FaaS-MALD}, it does
not optimize or certify a Lagrangian bound. Unlike \textsc{FaaS-MAPG}, its
stopping reasons do not certify a Nash equilibrium. FaaS-MAGCAA should
therefore be read as an implementation-faithful greedy allocation heuristic
whose empirical results can be compared under the same FRALB local-solve,
feasibility, evaluation, and persistence pipeline, not as a transfer of the
original GCAA theorem or as an alternative optimality certificate.
